\documentclass[11pt]{article}

% ---------- Packages ----------
\usepackage[margin=1in]{geometry}
\usepackage{amsmath}
\usepackage{amssymb}
\usepackage{mathtools}
\usepackage{graphicx}
\usepackage{booktabs}
\usepackage{hyperref}
\usepackage{xcolor}
\usepackage{listings}
\usepackage{enumitem}

% ---------- Hyperlink setup ----------
\hypersetup{
    colorlinks=true,
    linkcolor=blue,
    urlcolor=blue,
    citecolor=blue
}

% ---------- Code style ----------
\lstset{
    basicstyle=\ttfamily\small,
    frame=single,
    breaklines=true,
    columns=fullflexible
}

% ---------- Document Info ----------
\title{Assignment 1}
\author{Bill Hsu}
\date{\today}

\begin{document}

\maketitle

% =========================
\section{Introduction}
% =========================

This is a basic LaTeX template for homework and technical reports.

Inline math looks like this: $y = Wx + b$.

A standalone equation:

\[
y = Wx + b
\]

% =========================
\section{Equations}
% =========================

A numbered equation:

\begin{equation}
    E = mc^2
    \label{eq:energy}
\end{equation}

Refer to it using Equation~\ref{eq:energy}.

Multiple aligned equations:

\begin{align}
    Q &= XW_Q \\
    K &= XW_K \\
    V &= XW_V
\end{align}

Transformer attention:

\begin{equation}
    \mathrm{Attention}(Q,K,V)
    =
    \mathrm{softmax}
    \left(
    \frac{QK^\top}{\sqrt{d_k}}
    \right)V
\end{equation}

Common math symbols:
\begin{itemize}
    \item Fractions: $\frac{a}{b}$
    \item Square roots: $\sqrt{x}$
    \item Subscripts: $x_i$
    \item Superscripts: $x^2$
    \item Transpose: $A^\top$
    \item Bold vectors: $\mathbf{x}$
    \item Summation: $\sum_{i=1}^{n} x_i$
    \item Softmax function: $\mathrm{softmax}(z)$
\end{itemize}


% =========================
\section{Matrices and Vectors}
% =========================

A vector:

\[
\mathbf{x}
=
\begin{bmatrix}
x_1 \\
x_2 \\
x_3
\end{bmatrix}
\]

A matrix:

\[
A
=
\begin{bmatrix}
1 & 2 \\
3 & 4
\end{bmatrix}
\]

Useful notation:

\[
A^\top,
\qquad
A^{-1},
\qquad
\|\mathbf{x}\|_2,
\qquad
\sum_{i=1}^{n} x_i
\]

% =========================
\section{Lists}
% =========================

Bullet list:

\begin{itemize}
    \item First item
    \item Second item
    \item Third item
\end{itemize}

Numbered list:

\begin{enumerate}
    \item First step
    \item Second step
    \item Third step
\end{enumerate}

% =========================
\section{Figures}
% =========================

% \begin{figure}[ht]
%     \centering
%     \includegraphics[width=0.7\textwidth]{example.png}
%     \caption{Example figure.}
%     \label{fig:example}
% \end{figure}

% Refer to it using Figure~\ref{fig:example}.

% =========================
\section{Tables}
% =========================

\begin{table}[ht]
    \centering
    \begin{tabular}{lcc}
        \toprule
        Model & Accuracy & Loss \\
        \midrule
        Model A & 92.1\% & 0.31 \\
        Model B & 94.5\% & 0.24 \\
        \bottomrule
    \end{tabular}
    \caption{Example results.}
    \label{tab:results}
\end{table}

% =========================
\section{Code}
% =========================

\begin{lstlisting}[language=Python]
import torch

x = torch.randn(4, 8)
y = x @ x.T

print(y.shape)
\end{lstlisting}

% =========================
\section{Useful Formatting}
% =========================

\textbf{Bold text}

\textit{Italic text}

\texttt{monospace text}

A URL:

\url{https://www.example.com}

% =========================
\section{Conclusion}
% =========================

Write your conclusion here.

\end{document}